\subsection{RESTRICTIONS AND EXTENSIONS OF FUNCTIONS}

Two functions f and g are said to agree (or coincide) on a set E if E is contained in the domains of f and g and if \(f(x) = g(x)\) for all \(x \in E\) . Two functions which have the same graph agree on their domain. To say that f = g is to say that f and g have the same domain A and the same target B, and that they agree on A.

Let \(f = (\mathbf{F}, \mathbf{A}, \mathbf{B})\) and \(g = (\mathbf{G}, \mathbf{C}, \mathbf{D})\) be two functions. To say that \(\mathbf{F} \subset \mathbf{G}\) is to say that the domain \(\mathbf{A}\) of \(f\) is contained in the domain \(\mathbf{C}\) of \(g\) and that \(g\) agrees with \(f\) on \(\mathbf{A}\) . If also \(\mathbf{B} \subset \mathbf{D}\) , then \(g\) is said to be an extension of \(f\) (more precisely, an extension of \(f\) to \(\mathbf{C}\) ), and \(g\) is said to extend \(f\) (to \(\mathbf{C}\) ). When \(g\) is called a family of elements of \(\mathbf{D}\) , \(f\) is said to be a subfamily of \(g\) .

Let \(f\) be a function and let \(A\) be a subset of the domain of \(f\) . It is immediate that the relation `` \(x \in A\) and \(y = f(x)\) '' has a graph \(G\) with respect to \(x\) and \(y\) , that this graph is functional, and that \(A\) is its domain; the function whose graph is \(G\) , which has the same target as \(f\) , is called the restriction of f to A, and is sometimes denoted by \(f|A\) . A function is an extension of any of its restrictions. If two functions f, g have the same target and agree on a set E, then their restrictions to E are equal.

